% English version of paper/ja_v3/06_実験.md (ver3 draft 2, 2026-11-17). Ishii et al. 2022 form: method -> comparisons (what each numbered comparison verifies) -> results (subheadings name the comparisons; table reading, one paragraph per comparison, a summary paragraph with the answer in one sentence). Tables 8 and 9 of draft 1 are merged into Table 5; the comparator table moved to Table S3. Counts not in draft 1 ((D4) - (D1): 8 of 12, 4 above zero; (E2) - (E0): 10 of 12, 7 above zero) were counted from the values of Table 5.
\FloatBarrier
\section{Experiments: evaluation of the proposal}\label{sec:experiments}

\subsection{Procedure}\label{sec:procedure}
The evaluation targets are AFEW, IEMOCAP (improvised), MOSEI and MELD, with K = 5, 10 and 25 and 50 repetitions (5 folds $\times$ 10 draws). K labels, candidates, test utterances and product components are those of the (C1) evaluation run; only the classifiers measuring strengths and the AVT retrained on the proposal's pseudo-labels, (E2), are new. The proposal needs K labels, so there is no K = 0 condition. Intervals are those of the LOCO setting (Sect.~\ref{sec:setup}). CREMA-D (audio and visual) is also evaluated with 50 repetitions and four source corpora (Supplementary Information S7).

\subsection{Comparisons}\label{sec:compare}
Group D compares the estimation-time decisions of conditions that use the same product components and K labels and differ only in how predictions are combined; these comparisons do not depend on the retrained classifier. Group E compares the UAR of retrained AVT classifiers.
\begin{itemize}
\item \textbf{(D4) $-$ (D0), (D1):} the product and sum rules \cite{kittler1998} combine the three classifiers with equal weights and use no K labels. This verifies whether classifier-by-emotion weights raise UAR.
\item \textbf{(D4) $-$ (D2):} (D2) fits per-classifier inverse temperatures to the K-label log-likelihood (temperature scaling \cite{guo2017calibration}); like the proposal, it fits three quantities on the K labels.
\item \textbf{(D4) $-$ (D3):} (D3) selects the product component with the highest K-label strength and uses its decisions. This verifies whether combining the three classifiers itself raises UAR. (D3) was added after the evaluation run from its saved outputs (outside the evaluation protocol), with the same label-free intercept shift as the proposal.
\item \textbf{(D4) $-$ (D5), (D6):} (D5) sets $\beta=0$ and (D6) replaces the three strengths by their K-label mean. These verify the contributions of the per-emotion profiles and the strength differences.
\item \textbf{(E2) $-$ (E0), (E1):} whether the AVT retrained on the proposal's pseudo-labels has higher UAR than the AVT before retraining and the AVT retrained on the product's pseudo-labels.
\end{itemize}
Comparisons with three further comparators that fit weights on the K labels and with four further combinations of the same predictions are in Supplementary Information S6, and the content and references of all comparators are in Table~S3.

\subsection{Results}\label{sec:rq2-results}
Table~\ref{tab:rq2} gives the UAR of groups D and E and the differences from the proposal. For each corpus, the upper rows are estimation-time decisions and the lower rows retrained AVT classifiers.

\begin{table}[!htbp]
\caption{Evaluation of the proposal: UAR (points, test utterances, mean of 50 repetitions). (D4): the proposal's estimation-time decisions; (E2): the AVT retrained on its pseudo-labels. Parentheses: (D4) $-$ row for group D and (E2) $-$ row for group E, with 95\% intervals without correction; rows (D5) and (D6) give only the difference. * interval above zero (proposal higher); $\dagger$ below zero (proposal lower). (D3) was added after the evaluation run. Row names abbreviate those of Table~\ref{tab:conditions} (Mean: probability mean; Temp.: temperature calibration; Strongest: the single classifier strongest on the K labels; No profile: without per-emotion profiles; Equal str.: without strength differences; K-label: the AVT of (C0); Product: the AVT retrained on the product's pseudo-labels). CREMA-D (audio and visual) is in Table~S6}\label{tab:rq2}
\footnotesize\setlength{\tabcolsep}{1pt}
\begin{tabular}{@{}llll@{}}
\toprule
Row & K = 5 & K = 10 & K = 25 \\
\midrule
\multicolumn{4}{@{}l}{\textbf{AFEW}} \\
(D4) Proposal & 71.36 & 74.39 & 76.58 \\
(D0) Product & 69.23 (+2.13 [+1.0, +3.3]*) & 70.88 (+3.51 [+1.5, +5.5]*) & 73.45 (+3.14 [+1.9, +4.4]*) \\
(D1) Mean & 69.22 (+2.14 [+1.0, +3.2]*) & 71.02 (+3.38 [+1.9, +4.9]*) & 72.85 (+3.73 [+2.5, +5.0]*) \\
(D2) Temp. & 74.30 ($-$2.94 [$-$5.6, $-$0.3]$\dagger$) & 74.84 ($-$0.45 [$-$1.5, +0.6]) & 76.31 (+0.27 [$-$0.6, +1.1]) \\
(D3) Strongest & 68.65 (+2.70 [$-$3.3, +8.7]) & 70.21 (+4.18 [$-$1.2, +9.6]) & 72.79 (+3.79 [$-$1.5, +9.1]) \\
(D5) No profile & +0.38 [+0.03, +0.7]* & +0.45 [$-$0.1, +1.0] & +0.36 [$-$0.2, +1.0] \\
(D6) Equal str. & +1.74 [+0.4, +3.1]* & +3.20 [+1.1, +5.3]* & +2.57 [+1.4, +3.7]* \\
(E2) Proposal & 71.78 & 74.61 & 77.05 \\
(E0) K-label & 69.74 (+2.04 [+0.5, +3.6]*) & 71.56 (+3.05 [+1.8, +4.3]*) & 75.22 (+1.83 [+0.9, +2.8]*) \\
(E1) Product & 70.77 (+1.01 [+0.2, +1.8]*) & 72.66 (+1.95 [+0.5, +3.4]*) & 75.84 (+1.21 [+0.4, +2.0]*) \\
\midrule
\multicolumn{4}{@{}l}{\textbf{IEMOCAP (improvised)}} \\
(D4) Proposal & 72.56 & 73.80 & 75.03 \\
(D0) Product & 72.07 (+0.49 [$-$0.2, +1.2]) & 73.37 (+0.43 [$-$0.2, +1.1]) & 74.53 (+0.50 [+0.01, +1.0]*) \\
(D1) Mean & 73.03 ($-$0.46 [$-$2.8, +1.9]) & 74.25 ($-$0.46 [$-$3.1, +2.2]) & 75.52 ($-$0.49 [$-$3.2, +2.2]) \\
(D2) Temp. & 71.69 (+0.87 [+0.5, +1.2]*) & 73.16 (+0.64 [$-$0.1, +1.4]) & 74.37 (+0.67 [$-$0.2, +1.6]) \\
(D3) Strongest & 67.05 (+5.52 [+1.3, +9.8]*) & 68.42 (+5.38 [$-$0.5, +11.2]) & 71.51 (+3.52 [$-$3.6, +10.7]) \\
(D5) No profile & +0.24 [+0.03, +0.5]* & +0.23 [$-$0.05, +0.5] & +0.20 [$-$0.1, +0.5] \\
(D6) Equal str. & +0.38 [$-$0.2, +1.0] & +0.06 [$-$0.6, +0.7] & +0.22 [$-$0.3, +0.7] \\
(E2) Proposal & 72.81 & 73.86 & 74.95 \\
(E0) K-label & 72.55 (+0.25 [$-$4.5, +5.0]) & 73.49 (+0.36 [$-$3.7, +4.4]) & 75.51 ($-$0.56 [$-$4.7, +3.6]) \\
(E1) Product & 72.71 (+0.10 [$-$0.6, +0.8]) & 73.73 (+0.13 [$-$0.3, +0.6]) & 74.68 (+0.27 [$-$0.1, +0.6]) \\
\midrule
\multicolumn{4}{@{}l}{\textbf{MOSEI}} \\
(D4) Proposal & 41.16 & 42.82 & 44.47 \\
(D0) Product & 40.97 (+0.19 [$-$0.2, +0.6]) & 42.53 (+0.29 [$-$0.03, +0.6]) & 44.07 (+0.40 [+0.1, +0.7]*) \\
(D1) Mean & 41.01 (+0.15 [$-$0.2, +0.5]) & 42.56 (+0.26 [$-$0.1, +0.6]) & 43.94 (+0.52 [+0.2, +0.9]*) \\
(D2) Temp. & 39.98 (+1.18 [+0.7, +1.6]*) & 42.14 (+0.69 [+0.3, +1.0]*) & 44.11 (+0.36 [+0.1, +0.6]*) \\
(D3) Strongest & 36.83 (+4.33 [+3.5, +5.2]*) & 37.92 (+4.90 [+4.1, +5.7]*) & 39.20 (+5.27 [+4.5, +6.0]*) \\
(D5) No profile & +0.73 [+0.5, +1.0]* & +0.62 [+0.4, +0.9]* & +0.51 [+0.3, +0.7]* \\
(D6) Equal str. & $-$0.42 [$-$0.7, $-$0.1]$\dagger$ & $-$0.35 [$-$0.6, $-$0.1]$\dagger$ & $-$0.07 [$-$0.2, +0.1] \\
(E2) Proposal & 41.20 & 42.84 & 44.44 \\
(E0) K-label & 41.46 ($-$0.26 [$-$0.8, +0.3]) & 42.39 (+0.45 [$-$0.2, +1.0]) & 43.35 (+1.09 [+0.5, +1.7]*) \\
(E1) Product & 41.10 (+0.10 [$-$0.2, +0.4]) & 42.56 (+0.28 [$-$0.01, +0.6]) & 44.09 (+0.35 [+0.1, +0.6]*) \\
\midrule
\multicolumn{4}{@{}l}{\textbf{MELD}} \\
(D4) Proposal & 50.05 & 50.68 & 52.95 \\
(D0) Product & 50.15 ($-$0.09 [$-$0.6, +0.4]) & 50.20 (+0.48 [+0.01, +1.0]*) & 52.12 (+0.83 [+0.2, +1.5]*) \\
(D1) Mean & 50.28 ($-$0.22 [$-$1.4, +0.9]) & 50.38 (+0.30 [$-$0.7, +1.3]) & 52.30 (+0.65 [$-$0.1, +1.4]) \\
(D2) Temp. & 50.85 ($-$0.80 [$-$1.7, +0.1]) & 51.24 ($-$0.56 [$-$1.5, +0.3]) & 52.91 (+0.04 [$-$0.5, +0.6]) \\
(D3) Strongest & 46.79 (+3.27 [+1.6, +4.9]*) & 47.39 (+3.29 [+1.5, +5.1]*) & 50.08 (+2.88 [+1.4, +4.4]*) \\
(D5) No profile & +0.10 [$-$0.4, +0.6] & +0.14 [$-$0.5, +0.7] & +0.15 [$-$0.1, +0.4] \\
(D6) Equal str. & $-$0.19 [$-$0.7, +0.4] & +0.40 [$-$0.05, +0.8] & +0.75 [$-$0.004, +1.5] \\
(E2) Proposal & 50.07 & 50.38 & 52.73 \\
(E0) K-label & 47.59 (+2.48 [+1.8, +3.2]*) & 48.26 (+2.12 [+1.4, +2.8]*) & 50.49 (+2.24 [+1.6, +2.9]*) \\
(E1) Product & 50.13 ($-$0.05 [$-$0.6, +0.5]) & 50.13 (+0.25 [$-$0.2, +0.7]) & 51.90 (+0.82 [+0.2, +1.5]*) \\
\botrule\end{tabular}
\end{table}

\subsubsection{(D4) $-$ (D0), (D1): equal-weight combination rules}\label{sec:rq2-rules}
The proposal (D4) minus the product (D0) was positive in 11 of 12 conditions in point estimates, with intervals excluding zero in seven (AFEW at all K, IEMOCAP and MOSEI at K = 25, MELD at K = 10 and 25). The largest differences were on AFEW (+2.13 to +3.51\pt). On AFEW, the visual product component had a UAR more than 17\pt\ above the audio and text components (K = 5: visual 70.70, audio 53.24, text 45.76), and the visual classifier alone was higher than the product at K = 5 and 10 in point estimates (70.70 vs 69.23 and 71.93 vs 70.88). In the nine conditions on IEMOCAP, MOSEI and MELD, (D4) $-$ (D0) ranged from $-$0.09 to +0.83\pt. The difference from the probability mean, (D4) $-$ (D1), was positive in 8 of 12 conditions in point estimates, with intervals above zero on AFEW at all K and on MOSEI at K = 25 and below zero in none.

\subsubsection{(D4) $-$ (D2): a comparator fitted on the K labels}\label{sec:rq2-temp}
The difference from per-classifier temperature calibration, (D4) $-$ (D2), was positive in 8 of 12 conditions in point estimates, and the proposal was lower with an interval excluding zero only on AFEW at K = 5 ($-$2.94 [$-$5.6, $-$0.3], 71.36 vs 74.30). Comparisons with seven further comparators are in Supplementary Information S6.

\subsubsection{(D4) $-$ (D3): selecting the strongest single classifier}\label{sec:rq2-single}
The proposal (D4) was higher than the decisions of the single classifier strongest on the K labels (D3) in all 12 conditions in point estimates (+2.70 to +5.52\pt), with intervals above zero in seven (IEMOCAP at K = 5, MOSEI and MELD at all K). (D3) on AFEW at K = 5 (68.65) was lower than the visual classifier alone (70.70) because it selected a non-visual classifier in 3 of 50 repetitions and its intercepts were shifted (69.25 before the shift); the shift lowered (D3) on AFEW by 0.44--0.65\pt, and without the shift for either method (D4) $-$ (D3) was +1.86 to +6.11\pt\ in all 12 conditions (intervals above zero in eight). By contrast, per-emotion recall of (D4) was lower than that of (D3) in point estimates for Happy on AFEW ($-$6.42 to $-$4.39\pt), Neutral on MELD ($-$4.32 to $-$2.74\pt) and Sad on IEMOCAP at K = 25 ($-$0.21\pt; no per-emotion intervals computed). Splitting this difference, (D0) $-$ (D3) (+2.04 to +5.02\pt; intervals above zero in 7 of 9) accounted for most of it on IEMOCAP, MOSEI and MELD, and (D4) $-$ (D0) on AFEW ((D0) $-$ (D3) +0.57 to +0.67\pt), suggesting that multiplying the three classifiers' predictions raised UAR over selecting one on IEMOCAP, MOSEI and MELD, and changing their weights did so on AFEW.

\subsubsection{(D4) $-$ (D5), (D6): versions with one part removed}\label{sec:rq2-ablation}
On MOSEI, the difference from the proposal without per-emotion profiles, (D4) $-$ (D5), had intervals above zero at all K (+0.51 to +0.73\pt), and so did CREMA-D (Supplementary Information S7). On AFEW, the difference from the proposal without strength differences, (D4) $-$ (D6), had intervals above zero at all K (+1.74 to +3.20\pt), accounting for most of AFEW's (D4) $-$ (D0) (+2.13 to +3.51\pt). Counted after the evaluation over its 50 repetitions, the classifier with the highest K-label strength was the classifier with the highest strength on all candidate labels (visual on AFEW, audio on IEMOCAP, text on MELD) in 0.70--1.00 of repetitions, whereas the classifier with the highest source-measured strength was a different classifier on all three corpora. On MOSEI, where the top two strengths were close, this fraction was 0.30--0.38 (values per corpus and K in Supplementary Information S5).

\subsubsection{(E2) $-$ (E0), (E1): retrained AVT}\label{sec:rq2-retrained}
The AVT retrained on the proposal's pseudo-labels (E2) had higher UAR than the AVT retrained on the product's pseudo-labels (E1) in 11 of 12 conditions in point estimates, with intervals above zero in five (AFEW at all K, MOSEI and MELD at K = 25); on AFEW the differences were +1.01, +1.95 and +1.21\pt, and on the other three corpora $-$0.05 to +0.82\pt. The difference from the AVT before retraining, (E2) $-$ (E0), was positive in 10 of 12 conditions in point estimates, with intervals above zero in seven. On CREMA-D (audio and visual), the estimation-time (D4) $-$ (D0) and the retrained AV (E2) $-$ (E1) were +0.31 to +0.63\pt\ with all six intervals excluding zero (Supplementary Information S7).

\subsubsection{Summary of the results for RQ2}
The answer to RQ2 is therefore as follows. Combining predictions with reliabilities whose strengths come from the K labels and per-emotion profiles from the sources gave estimation-time UAR higher than the equal-weight product in 11 of 12 conditions in point estimates (by 2.1--3.5\pt\ on AFEW), and the AVT retrained on the resulting pseudo-labels had higher UAR than the AVT retrained on the product's pseudo-labels in 11 of 12. Differences from temperature calibration on the K labels, (D2), varied in sign.
